\cleardoublepage

\chapter*{Resumen\markboth{Resumen}{Resumen}}

La asistencia sanitaria domiciliaria constituye uno de los principales retos asociados al envejecimiento de la población y al aumento de las enfermedades crónicas. En este contexto, la correcta administración de la medicación y el seguimiento periódico de determinados parámetros fisiológicos resultan fundamentales para mejorar la calidad de vida de los pacientes. El avance de las tecnologías embebidas y de los dispositivos móviles ha permitido el desarrollo de nuevas herramientas destinadas a facilitar estas tareas en entornos domésticos.

A pesar de estos avances, muchas de las soluciones existentes presentan costes elevados o funcionalidades limitadas. Además, la gestión de la medicación y la monitorización de parámetros biomédicos suelen encontrarse separadas en diferentes dispositivos, dificultando la supervisión conjunta de la información. Por ello, resulta de interés desarrollar soluciones accesibles que permitan integrar estas funcionalidades en una única plataforma.

Para abordar esta problemática, en el presente trabajo se ha desarrollado MyMeds, un sistema inteligente de asistencia sanitaria domiciliaria compuesto por un dispositivo físico basado en un microcontrolador ESP32 con pantalla táctil y una aplicación móvil Android. El sistema permite gestionar medicamentos y tomas programadas, sincronizar información mediante WiFi, almacenar historiales de medicación y monitorizar la frecuencia cardíaca mediante un sensor biomédico. Asimismo, incorpora mecanismos de almacenamiento persistente y autenticación mediante PIN para garantizar la conservación y protección de la información.

Con el objetivo de validar la solución propuesta, se han realizado diferentes experimentos relacionados con la monitorización biomédica, la sincronización de datos, la persistencia de la información y la seguridad de las comunicaciones. Los resultados obtenidos muestran un funcionamiento satisfactorio de las funcionalidades implementadas y demuestran la viabilidad técnica del sistema desarrollado como herramienta de apoyo para la gestión de tratamientos farmacológicos en entornos domésticos.

